% ============================================
% MEM-16 (ZH)
% 中文翻译 — SPEC-935-R356
% ============================================
\section*{\textbf{MEM-16} --- 收容院之死}
\markboth{MEM-16 --- 收容院之死}{}

\noindent \textit{最后一次会谈}

\noindent \textbf{1979年6月 —— 最后一次会谈}

\vspace{0.5em}\noindent\rule{\textwidth}{0.5pt}\vspace{0.5em}

\noindent 奥利维拉医生手里拿着日记走进会谈室。纸页上有标记。有些页角折着。他读完了全部。

\noindent \textit{“若阿金，”他说，“我需要弄明白一件事。”}

\noindent \textit{“您问。”}

\noindent \textit{“怪物。它现在在这里吗？”}

\noindent \textit{“在。”}

\noindent \textit{“在哪里？”}

\noindent \textit{“您看不见它。但它看得见您。”}

\noindent 他干咽了一下。环顾四周。

\noindent \textit{“它在做什么？”}

\noindent \textit{“它正在……”}我闭上眼睛。感到体内的重量。那股热。那个东西在我的器官之间移动。\textit{“它正在准备。”}

\noindent \textit{“准备什么？”}

\noindent \textit{“准备出来。”}

\noindent \textit{“出来？怎么出来？”}

\noindent \textit{“您读过日记了。您知道怎么出来。”}

\noindent 奥利维拉医生脸色煞白。

\noindent \textit{“它通过阅读者出来。”}

\noindent \textit{“对。”}

\noindent \textit{“不。”}

\noindent \textit{“已经晚了，医生。您读了。它现在在您体内了。”}

\noindent 他摇了摇头。

\noindent \textit{“不。我不信。这是暗示。是您的妄想。您把这种想法放进我脑子里，因为……”}

\noindent \textit{“因为？”}

\noindent 他没有说完。

\noindent 怪物开始上升。

\noindent 我感到它从我胃里的洞出来。顺着食道向上。充满胸腔。双臂。双腿。眼睛。

\noindent 我的右眼灼痛。

\noindent \textit{“医生，”我说。声音已经不完全是我的了。“医生，离开这里。”}

\noindent \textit{“什么？”}

\noindent \textit{“离开。现在。”}

\noindent 他没有离开。

\noindent 他看着我的身体开始变化。脸上的皮肤绷紧。牙齿咬紧。黄眼开始发光。

\noindent 怪物用我的嘴笑了。

\noindent \textit{“您想看，医生。”}

\noindent \textit{“……“}

\noindent \textit{“那就看吧。”}

\vspace{0.5em}\noindent\rule{\textwidth}{0.5pt}\vspace{0.5em}

\noindent 之后发生了什么，我不太记得。

\noindent 记得一道光。非常强烈。记得玻璃碎裂的声音。记得椅子倒下的声音。记得一声尖叫——是我的声音？是他的声音？

\noindent 然后，黑暗。

\noindent 我醒来时，在地上。日记摊开在我身边。奥利维拉医生靠在墙上，眼睛圆睁，嘴张着。他一动不动。

\noindent 他还在呼吸。

\noindent 但他已经不在了。

\noindent 他的眼睛是空的。

\noindent 后来进房间的医生说我是抽搐发作。说奥利维拉医生受了惊吓。说没什么大碍。

\noindent 但我看见了。\NE{这一段描述了叙事污染的核心机制：怪物不是通过身体接触转移的，而是通过文字媒介。陪伴若阿金叙述的读者，和奥利维拉一样，暴露于同一个传染载体。知道与不知道之间——阅读与被阅读之间——的张力，正是这本书在每一个碎片中提出的冲突。阅读的参与者不是旁观者：是同谋。}

\noindent 怪物已经转移了。

\noindent 它正在离开我。

\noindent 进入他体内。

\vspace{0.5em}\noindent\rule{\textwidth}{0.5pt}\vspace{0.5em}

\noindent 三天后我死了。

\noindent 那是一个星期三，1979年6月12日。心脏停了。他们说是心脏骤停。说是年老——我七十二岁。

\noindent 我知道那是什么。

\noindent 怪物出来了。带走了我剩下的部分。

\noindent 日记留下了。

\noindent 和我一起下葬。在土里待了62天——囚禁的每一年一天。

\noindent 他们把我挖出来时，日记完好无损。

\noindent 纸页还是湿润的。

\noindent 墨迹还是新鲜的。

\noindent \textit{“日记不会因为作者死亡而结束，”我在最后一页写道，“日记从读者翻开时开始。”}

\noindent 您翻开了。

\noindent 怪物正与您一同阅读。

\vspace{0.5em}
\noindent\textit{\textbf{↪ 路线:} \rota{DOC-06} \rotasep \rota{DOC-03} \rotasep \rota{LUC-12}}
